% =============================================================================
% Bloque 6: Metodología (slides 12–16)
% =============================================================================

\section{Metodología}

% --- Slide 12: Pipeline arquitectónico ---
\begin{frame}{Pipeline: del texto real al clasificador entrenado}
\centering
\begin{tikzpicture}[
    every node/.style={font=\scriptsize},
    box/.style={draw=uaiblue, rounded corners, fill=uaibluelight!15,
                minimum width=2cm, minimum height=0.9cm, text width=2cm, align=center},
    box2/.style={draw=uaiorange, very thick, rounded corners, fill=uaiorange!15,
                 minimum width=2cm, minimum height=0.9cm, text width=2cm, align=center},
    box3/.style={draw=uaigreen, very thick, rounded corners, fill=uaigreen!15,
                 minimum width=2cm, minimum height=0.9cm, text width=2cm, align=center},
    arrow/.style={->, >=Stealth, thick, uaiblue!70}
  ]
  % Fila superior
  \node[box]  (real)  at (0,2)    {\textbf{Datos\\reales}\\\tiny 10/25/50 por clase};
  \node[box]  (enc)   at (2.6,2)  {\textbf{Al mapa}\\\tiny embeddings};
  \node[box]  (llm)   at (5.2,2)  {\textbf{IA genera}\\\tiny Gemini\\\tiny $3\times$ de sobra};
  \node[box]  (cands) at (7.8,2)  {\textbf{Candidatas}\\\tiny puntos en el mapa};
  \node[box2] (filt)  at (10.4,2) {\textbf{Control de calidad}\\\tiny filtro geométrico};

  % Fila inferior
  \node[box2] (weight) at (10.4,0)  {\textbf{Pesos}\\\tiny según calidad};
  \node[box3] (clf)    at (5.2,0)   {\textbf{Clasificador}\\\tiny 5 tipos distintos};
  \node[box]  (eval)   at (0,0)     {\textbf{Evaluación}\\\tiny con datos nunca vistos};

  % Flechas
  \draw[arrow] (real)  -- (enc);
  \draw[arrow] (enc)   -- (llm);
  \draw[arrow] (llm)   -- (cands);
  \draw[arrow] (cands) -- (filt);
  \draw[arrow] (filt)  -- (weight);
  \draw[arrow] (weight) -- (clf);
  \draw[arrow] (real.south)  .. controls +(0,-0.6) and +(-0.5,0) .. (clf.west);
  \draw[arrow] (clf)   -- (eval);
  % Atajo punteado: filtro binario (sin ponderación suave)
  \draw[->, >=Stealth, thick, uaibluedark!60, dashed]
       (filt.south west) .. controls +(-1.2,-0.7) and +(1.0,0.7) ..
       node[midway, above, font=\tiny\itshape, color=uaibluedark!70] {sin pesos}
       (clf.north east);
\end{tikzpicture}

\vspace{0.5em}
\mensaje{Las dos piezas nuevas (control de calidad y pesos) no tocan la IA: funcionan con cualquiera.}
\end{frame}

% --- Slide 13: Protocolo experimental (vuelve, antes de los datasets) ---
\begin{frame}{Protocolo experimental: 8.000+ experimentos}
\small
\begin{columns}[T]

\begin{column}{0.55\textwidth}
\textbf{Probamos la idea en todas las combinaciones:}
\begin{itemize}\setlength\itemsep{0pt}
  \item \textbf{13 conjuntos de datos} en 9 dominios.
  \item \textbf{Escasez:} 10 / 25 / 50 ejemplos por clase.
  \item \textbf{5 tipos} de clasificador.
  \item \textbf{10 métodos} de aumentación en competencia.
  \item \textbf{5 IAs generadoras} (Gemini, GPT-5, Claude, Kimi, GLM).
  \item \textbf{4 mapas} (modelos de embeddings) distintos.
  \item \textbf{5 repeticiones} de cada experimento.
\end{itemize}
\end{column}

\begin{column}{0.45\textwidth}
\textbf{Cómo comparamos:}
\begin{itemize}\setlength\itemsep{0pt}
  \item Nota de calidad: \textbf{macro F1} (siguiente slide).
  \item Cada método compite contra \textbf{SMOTE} en el mismo escenario exacto.
  \item Pruebas estadísticas estrictas: descartan que la diferencia sea \alert{azar}.
  \item También contamos \textbf{cuántas veces gana} cada método.
\end{itemize}

\vspace{0.3em}
\setbeamercolor{block title}{fg=white,bg=uaiorange}
\begin{block}{Volumen}
\centering\bfseries\large 8.000+ experimentos
\end{block}
\end{column}

\end{columns}
\end{frame}

% --- Slide 13b: Métrica de evaluación — Macro F1 (sin fórmulas, versión divulgación) ---
\begin{frame}{¿Cómo medimos el éxito? Una nota justa: macro F1}
\begin{columns}[T]

\begin{column}{0.55\textwidth}
\small
\textbf{El acierto global engaña:} un modelo que responde siempre «clase común» acierta el 95\%\dots\ ignorando por completo a la clase rara.

\vspace{0.5em}
\textbf{Macro F1} hace algo más justo:
\begin{enumerate}\setlength\itemsep{3pt}
  \item Le pone una \textbf{nota de 0 a 100 a cada clase}: qué tan bien la detecta, sin confundirla con otras.
  \item \textbf{Promedia las notas}: cada clase pesa lo mismo, tenga 16 ejemplos o 720.
\end{enumerate}

\vspace{0.5em}
\begin{insight}
\scriptsize
Como el promedio de notas por ramo: sacarse un 7 en el ramo fácil \textbf{no tapa} el 2 en el difícil.
\end{insight}
\end{column}

\begin{column}{0.45\textwidth}
\setbeamercolor{block title}{fg=white,bg=uaiblue}
\begin{block}{\small Por qué nos sirve}
\scriptsize
\begin{itemize}\setlength\itemsep{1pt}
  \item Un error en la clase rara \textbf{duele igual} que uno en la clase grande.
  \item Si el modelo \emph{ignora} una clase chica, la nota cae de inmediato.
\end{itemize}
\end{block}

\vspace{0.4em}
\setbeamercolor{block title}{fg=white,bg=uaiorange}
\begin{block}{\small Cómo leerán los resultados}
\scriptsize
\begin{itemize}\setlength\itemsep{1pt}
  \item $\Delta$F1 $=$ cuántos \textbf{puntos de nota} sube (o baja) un método frente a SMOTE.
  \item «pp» $=$ puntos porcentuales: $+2$ pp $=$ subir de 71 a 73.
  \item Siempre evaluado sobre datos que el modelo \textbf{nunca vio}.
\end{itemize}
\end{block}
\end{column}

\end{columns}

\mensaje{Macro F1: la métrica honesta cuando hay clases raras.}
\end{frame}

% --- Slide 14: Conjuntos de datos (tabla) ---
\begin{frame}{Conjuntos de datos: 13 corpus en 9 dominios}
\centering
\scriptsize
\renewcommand{\arraystretch}{0.92}
\begin{tabular}{@{}llccr@{}}
\toprule
\textbf{Conjunto} & \textbf{Dominio} & \textbf{Clases} & \textbf{$n$-shots} & \textbf{Test} \\
\midrule
\rowcolor{uaibluelight!12}\multicolumn{5}{l}{\textit{\color{uaibluedark}Clasificación de texto --- experimento principal}} \\
\dataset{sms\_spam}              & Detección de spam        & 2  & 10, 25, 50 & 1.394 \\
\dataset{hate\_speech\_davidson} & Discurso de odio         & 3  & 10, 25, 50 & 4.948 \\
\dataset{20newsgroups}           & Noticias (4 cat.)        & 4  & 10, 25, 50 & 7.532 \\
\dataset{ag\_news}               & Noticias (4 cat.)        & 4  & 10, 25, 50 & 7.600 \\
\dataset{emotion}                & Emociones                & 6  & 10, 25, 50 & 2.000 \\
\dataset{dbpedia14}              & Clasificación de temas   & 14 & 10, 25, 50 & 5.000 \\
\dataset{20newsgroups\_20class}  & Noticias (20 cat.)       & 20 & 10, 25, 50 & 7.532 \\
\midrule
\rowcolor{uaiorange!12}\multicolumn{5}{l}{\textit{\color{uaibluedark}Adicionales --- escalabilidad por número de clases}} \\
\dataset{trec6}                  & Preguntas                & 6   & 10, 25, 50 & 500 \\
\dataset{banking77}              & Intenciones bancarias    & 77  & 10, 25, 50 & 3.080 \\
\dataset{clinc150}               & Intenciones de diálogo   & 150 & 10, 25, 50 & 4.500 \\
\midrule
\rowcolor{uaigreen!12}\multicolumn{5}{l}{\textit{\color{uaibluedark}Reconocimiento de entidades nombradas (NER)}} \\
\dataset{MultiNERD}              & Entidades generales      & 5 tipos & 10, 25, 50 & 10.000 \\
\dataset{WikiANN}                & Entidades Wikipedia      & 3 tipos & 10, 25, 50 & 10.000 \\
\dataset{Few-NERD}               & Entidades detalladas     & 8 tipos & 10, 25, 50 & 10.000 \\
\bottomrule
\end{tabular}

\vspace{0.2em}
{\scriptsize\itshape $n$-shot estricto: 10 / 25 / 50 ejemplos por clase. Rango de complejidad: de binario (\dataset{sms\_spam}) a 150 clases (\dataset{clinc150}).}
\end{frame}

% --- Slide 16: ¿Qué es NER y por qué la usamos? ---
\begin{frame}{¿Qué es NER y por qué la usamos?}

\setbeamercolor{block title}{fg=white,bg=uaibluedark}
\begin{block}{\small NER (\textit{Named Entity Recognition})}
\scriptsize
Tarea de NLP que \textbf{detecta nombres propios dentro del texto y los clasifica por tipo}: personas, lugares, organizaciones, fechas, etc.

\vspace{0.3em}
\textit{Ejemplo:} «Llamé a \colorbox{uaiblue!25}{\strut\,María\,} en \colorbox{uaiorange!25}{\strut\,Madrid\,} el \colorbox{uaigreen!25}{\strut\,2 de marzo\,}.»\hspace{0.5em}{\tiny \textcolor{uaiblue}{$\uparrow$PER}\hspace{1em}\textcolor{uaiorange}{$\uparrow$LOC}\hspace{1em}\textcolor{uaigreen}{$\uparrow$DATE}}
\end{block}

\vspace{0.4em}
\begin{columns}[T]

\begin{column}{0.55\textwidth}
\small
\textbf{¿Por qué la usamos como prueba?}
\begin{itemize}\setlength\itemsep{1pt}\scriptsize
  \item Construimos el filtro para \textbf{clasificación}.
  \item NER es una tarea NLP \textbf{distinta}: en vez de etiquetar la oración entera, se etiquetan palabras dentro.
  \item Si el mismo filtro funciona en NER \textbf{sin cambiarlo} $\Rightarrow$ el filtro \alert{no depende de la tarea}.
\end{itemize}

\vspace{0.2em}
{\scriptsize Es la prueba más fuerte de que el método \textbf{generaliza}.}
\end{column}

\begin{column}{0.45\textwidth}
\setbeamercolor{block title}{fg=white,bg=uaigreen}
\begin{block}{\small 3 corpus, distinta complejidad}
\scriptsize
\begin{itemize}\setlength\itemsep{1pt}
  \item \dataset{WikiANN} --- 3 tipos (PER, ORG, LOC). \textbf{Fácil}.
  \item \dataset{MultiNERD} --- 5 tipos (+ DATE, MISC). Medio.
  \item \dataset{Few-NERD} --- 8 tipos jerárquicos. \textbf{Difícil}.
\end{itemize}
\end{block}
\end{column}

\end{columns}

\mensaje{Mismo filtro, otra tarea: si funciona acá, sirve en general.}
\end{frame}

% --- Slide 17: Los 4 filtros (versión divulgación) ---
\begin{frame}{Cuatro maneras de hacer el control de calidad}
\small
Probamos cuatro «porteros» distintos, del más permisivo al más estricto:

\vspace{0.5em}
\begin{tabular}{@{}p{2.9cm}p{9.2cm}@{}}
\toprule
\textbf{Filtro} & \textbf{La pregunta que le hace a cada candidata} \\
\midrule
\textit{Ninguno} (control)  & Deja pasar todo. Sirve de punto de comparación. \\[4pt]
\textit{LOF} (densidad)     & ¿Caíste en una zona \textbf{poblada} de tu clase, o en un descampado? \\[4pt]
\textit{Cascada} (niveles 1--4) & ¿Qué tan \textbf{cerca} caíste de los ejemplos reales? Ordena a todas las candidatas y se queda con las mejores. \\[4pt]
\textit{Combinado}          & Exige \textbf{dos condiciones a la vez} (densidad \textbf{y} parecido). El portero más estricto. \\
\bottomrule
\end{tabular}

\mensaje{Mismo mapa, cuatro criterios: ¿cuál elige mejor?}
\end{frame}

% --- Slide 14: Cascade — pieza central (versión divulgación, sin fórmulas) ---
\begin{frame}{La cascada: cuatro preguntas sobre cada candidata}
\begin{columns}[T]

\begin{column}{0.55\textwidth}
\small
Cada texto que escribe la IA recibe un \textbf{puntaje de calidad} según dónde cayó su punto en el mapa. Cuatro preguntas, de la más simple a la más elaborada:

\vspace{0.4em}
\begin{description}\setlength\itemsep{4pt}\small
  \item[Nivel 1 \alert{$\star$}] \textbf{Cercanía:} ¿qué tan cerca caíste del ejemplo real más próximo?
  \item[Nivel 2] \textbf{Parecido:} ¿te «pareces» a ese ejemplo? (similitud de significado)
  \item[Nivel 3] \textbf{Vecindario:} ¿tus vecinos en el mapa son de tu misma clase?
  \item[Nivel 4] \textbf{Centro:} ¿qué tan cerca estás del centro de tu clase?
\end{description}

\vspace{0.4em}
Las preguntas activas se combinan en un solo puntaje, y nos quedamos con las \textbf{$N$ mejores} candidatas del ranking.
\end{column}

\begin{column}{0.45\textwidth}
\centering
\begin{tikzpicture}[every node/.style={font=\tiny}]
  % Pirámide de niveles
  \fill[uaiorange!40] (-2.5,2.4) rectangle (2.5,3.0);
  \node[font=\scriptsize\bfseries] at (0,2.7) {Nivel 1: Cercanía $\star$};

  \fill[uaibluelight!30] (-2.2,1.7) rectangle (2.2,2.3);
  \node[font=\scriptsize] at (0,2.0) {Nivel 2: Parecido};

  \fill[uaibluelight!20] (-1.9,1.0) rectangle (1.9,1.6);
  \node[font=\scriptsize] at (0,1.3) {Nivel 3: Vecindario};

  \fill[uaibluelight!15] (-1.6,0.3) rectangle (1.6,0.9);
  \node[font=\scriptsize] at (0,0.6) {Nivel 4: Centro de clase};

  % Flecha de score
  \draw[->, very thick, uaiorange] (-3.2,0.3) -- (-3.2,3.0);
  \node[font=\tiny, rotate=90] at (-3.6,1.65) {puntaje $\uparrow$};
\end{tikzpicture}\\[4pt]
{\scriptsize Se queda con las \textbf{$N$ mejores} del ranking\\ (no hay nota mínima de corte).}

\vspace{0.5em}
\begin{insight}
\scriptsize \textbf{Adelanto:} la pregunta más simple —el Nivel 1, pura cercanía— resultará ser la única necesaria.
\end{insight}
\end{column}

\end{columns}
\end{frame}

% --- Slide 15: Ponderación suave (versión divulgación, sin ecuaciones) ---
\begin{frame}{Ponderación suave: no todas las muestras valen lo mismo}
\begin{columns}[T]

\begin{column}{0.5\textwidth}
\small
La segunda contribución: en vez de \textbf{aceptar o rechazar} en seco, cada muestra entra al entrenamiento con un \textbf{peso según su puntaje}.

\vspace{0.4em}
\begin{itemize}\setlength\itemsep{2pt}
  \item \textbf{Ejemplo real} $\Rightarrow$ peso máximo (1): en ellos confiamos siempre.
  \item \textbf{Candidata bien ubicada} $\Rightarrow$ peso alto.
  \item \textbf{Candidata dudosa} $\Rightarrow$ peso bajo: se usa, pero \alert{se le cree menos}.
\end{itemize}

\vspace{0.4em}
\begin{insight}
\scriptsize
Como un profesor que escucha a todo el curso, pero pondera más las respuestas de quienes muestran su desarrollo.
\end{insight}
\end{column}

\begin{column}{0.5\textwidth}
\centering
\begin{tikzpicture}[every node/.style={font=\tiny}]
  % Eje
  \draw[->,uaigray] (0,0) -- (4.5,0) node[right]{puntaje};
  \draw[->,uaigray] (0,0) -- (0,3.2) node[above]{peso};
  \node[font=\tiny] at (4,-0.25) {1};
  \node[font=\tiny] at (-0.2,3) {1};

  % Curva tau=0.5 (peaked) — w = s^2
  \draw[very thick, uaiorange, domain=0:4, samples=40, smooth] plot (\x, {3*(\x/4)^2});
  \node[uaiorange, font=\scriptsize, align=center] at (3.1,2.75) {curva exigente\\(la que usamos)};

  % Curva tau=1 (linear)
  \draw[thick, uaiblue, dashed, domain=0:4, samples=40] plot (\x, {3*(\x/4)});
  \node[uaiblue, font=\scriptsize, align=center] at (3.7,1.35) {curva\\neutra};

  \node[align=center, font=\tiny] at (2.2,-0.7) {Puntaje de calidad de la muestra};
\end{tikzpicture}\\[4pt]
{\scriptsize La curva exigente castiga fuerte a las dudosas\\y premia a las muy bien ubicadas.}
\end{column}

\end{columns}

\mensaje{El clasificador escucha a todas las muestras, pero les cree en proporción a su calidad.}
\end{frame}

